\subsection{Multicol -- Multiple Columns of text}

The \texttt{multicol} package divides part of the document into
multiple columns. The number of columns is specified when opening
the environment:

\begin{lstlisting}[style=latex]
\begin{multicols*}{4}

    Content of the first column

    \columnbreak

    Content of the second column

    \columnbreak

    Content of the third column

    \columnbreak

    Content of the fourth column

\end{multicols*}
\end{lstlisting}

The example above creates four columns. The
\verb|\columnbreak| command can be used to manually move to the next
column.

This is particularly useful when several short sections must be
arranged compactly on the same page:

\begin{lstlisting}[style=latex]
\begin{multicols*}{3}
    \input{sections/hooke}
    \input{sections/bending}

    \columnbreak

    \input{sections/torsion}
    \input{sections/power}

    \columnbreak

    \input{sections/buckling}
\end{multicols*}
\end{lstlisting}

The starred environment \texttt{multicols*} does not automatically
balance the height of the final columns. This is useful when the
position of the content is controlled manually using
\verb|\columnbreak|.